%% Rewritten Results section
%% All numbers verified against: run_full_analysis.R, run_ideas.R, fs_checks.R
%% Tables referenced: tab:r1_infra (R1_infra.tex), tab:r1_welfare (R1_welfare.tex),
%%   tab:r1_lights (R1_lights.tex), tab:r2a (R2A_sign_contrast.tex),
%%   tab:ideas_summary (ideas_summary.tex)

\section{Results}
\label{sec:results}

The results follow the logic of the theoretical framework. The average contrast between
technical and social-science officers is attenuated toward zero, exactly as
Proposition~\ref{prop:atten} anticipates when educational category proxies noisily for
capacity. Correcting for this attenuation by instrumenting recovers an effect of the
predicted sign and magnitude. The falsification outcome behaves as
Proposition~\ref{prop:mix} requires. The one prediction concerning the tenure profile
of the effect is not supported.

%% ------------------------------------------------------------------
\subsection{The Average Effect Is Attenuated}
\label{sec:results-ols}

Before turning to the joint test of comparative advantage it is useful to examine the
individual outcome regressions that underlie the two programme margins of interest.
Tables~\ref{tab:r1_infra} and~\ref{tab:r1_welfare} report engineer coefficients across
four successively demanding fixed-effect specifications---district only; district plus
year; district plus additive state and year; district plus state-by-year---for the
infrastructure and welfare outcomes respectively.

For the \emph{infrastructure} outcome (road completion rate), the engineer coefficient
is $+0.034$ under district fixed effects alone, shrinks to $+0.022$ once year fixed
effects absorb aggregate programme rollout, and remains stable at $+0.022$ when
additive state effects are added. Under the most saturated specification, state-by-year
fixed effects, the coefficient is $+0.018$. Across all four columns the sign is positive
and the point estimate is economically modest, lying between 1.8 and 3.4 percentage
points of the completion rate, but with standard errors of similar magnitude it is
statistically indistinguishable from zero throughout. Crucially, the stability of the
coefficient as controls are added argues against the confound story: the estimate does
not fall toward zero as state-specific trends are absorbed, which would be the signature
of positive assignment of engineer Collectors to high-investment districts. Instead the
coefficient is flat from specification (2) onward, suggesting the variation is
idiosyncratic rather than state-driven.

For the \emph{welfare} outcome (employment per capita), the engineer coefficient is
$+0.0007$ under district fixed effects, $+0.0012$ across specifications (2) and (3),
and $+0.0012$ under state-by-year fixed effects---a magnitude two orders of magnitude
smaller than the infrastructure point estimate, economically negligible and statistically
invisible throughout. The absence of a detectable welfare effect is consistent with
Proposition~\ref{prop:comp}: an engineer Collector who tilts effort toward infrastructure
should reduce welfare programme output, so a positive (or zero) welfare coefficient in
OLS is evidence of the attenuation that obscures the differential, not evidence of
the absence of a differential. The nightlights outcome (Table~\ref{tab:r1_lights})
displays the same near-zero pattern, with engineer coefficients of $-0.0007$ to $-0.0009$
across all specifications, a result taken up in Section~\ref{sec:results-falsification}.

The absence of detectable signal in any individual outcome does not require that the
mechanism fails; it requires only that the total effect on any single outcome be small
relative to its noise. Proposition~\ref{prop:atten} predicts this: the relevant signal
is the \emph{differential} across programme types, not the level for any one. To
recover that differential we pool the two programme outcomes into a single stacked
regression. Standardising roads and employment to comparable units and assigning
opposite signs---$Y_{\text{signed}} = z(\text{roads})$ for infrastructure,
$-z(\text{employment})$ for welfare---so that comparative advantage predicts a single
positive engineer coefficient, the OLS estimate is $+0.037$ ($\text{SE}=0.033$,
$p=0.26$, Table~\ref{tab:ideas_summary}, row~1). The estimate is of the correct sign
but statistically indistinguishable from zero.

An equivalent test, reported in Table~\ref{tab:r2a}, uses the explicit
interaction form: regressing the standardised outcome on engineer, the
engineer-by-infra-indicator interaction (whose coefficient equals
$\hat\beta_{\text{eng}}^{\text{infra}} - \hat\beta_{\text{eng}}^{\text{welfare}}$),
and the same fixed effects. The interaction coefficient is $+0.14$ ($\text{SE}=0.11$,
$p=0.21$) in the preferred additive-state specification and $+0.15$ ($\text{SE}=0.11$,
$p=0.18$) under state-by-year fixed effects. A Wald test of the null that the two
programme-specific engineer coefficients are equal yields $\chi^2(1)=0.84$ ($p=0.36$).
The direction is consistently correct and the magnitude---roughly $0.14$ standard
deviations of the outcome---is economically plausible for a posting-level intervention,
but the sample is insufficient to distinguish it from sampling noise.

An alternative operationalisation of the sign contrast constructs the log-ratio
$\log(\text{roads\_std} + c) - \log(\text{emp\_std} + c)$ directly, so that a positive
engineer coefficient signifies infrastructure output rising relative to welfare output
within the same district-term. Across three choices of the small constant $c$ the
engineer coefficient ranges from $+0.063$ to $+0.086$ ($p=0.29$ to $0.51$). All
estimates are of the predicted sign; none is significant. This alternative operationalisation
corroborates the direction of the effect while confirming that OLS cannot detect it at
this sample size.

This pattern of correctly-signed, insignificant OLS estimates is not evidence against
the mechanism. Proposition~\ref{prop:atten} predicts precisely this outcome: because the
estimator is built on educational category, which constitutes one component of capacity
among three, the coefficient is scaled toward zero by the ratio
$\operatorname{Var}(\bar{e}(\tau))/\operatorname{Var}(e)$. The average contrast is
exactly the quantity the theory predicts will be least visible.

%% ------------------------------------------------------------------
\subsection{Instrumenting Recovers the Predicted Effect}
\label{sec:results-iv}

If the OLS null reflects attenuation rather than the absence of an effect, an instrument
that isolates the education-driven component of officer assignment should recover a
larger estimate. This is the content of Proposition~\ref{prop:atten}: instrumental-variables
estimation corrects classical measurement error, so the IV coefficient should exceed the
OLS coefficient in magnitude.

Panel~B of Table~\ref{tab:ideas_summary} confirms this ordering. Instrumenting engineer
assignment with the queue-composition share $Q_{st}$, the two-stage least-squares
estimate of the pooled comparative-advantage effect is $+0.337$
($\text{SE}=0.202$, $p=0.097$), significant at the ten percent level. The estimate is
roughly nine times the OLS coefficient, the direction and approximate magnitude of
correction that attenuation predicts.

The first stage is strong. Because $Q_{st}$ varies at the state-year level, a
state-by-year fixed effect would mechanically absorb it; identification therefore
uses district-by-outcome, state-by-outcome, and year-by-outcome fixed effects, which
partial out state-specific and year-specific confounds for each programme class
separately without conflating the two.\footnote{A specification adding a common
state-by-year fixed effect, which absorbs the instrument, confirms the mechanical
absorption: the instrument coefficient is identically zero under that specification,
a useful sanity check.} The Kleibergen--Paap $F$-statistic---computed as the squared
$t$-statistic on $Q_{st}$ in the first stage, the appropriate measure for a
just-identified instrument---is $34.5$ in the preferred specification. It remains
between $30$ and $39$ across all specifications that preserve the exogeneity argument,
comfortably above conventional thresholds for instrument strength. The instrument's
power does not depend on any particular control: dropping the electoral-turnover
control alone leaves $F=29.4$; dropping the tenure control alone gives $F=35.5$;
dropping both yields $F=30.4$. Identification rests on the queue-composition variation
formalised in Section~\ref{sec:identification}: the share of technical officers in a
state's active posting pool shifts the probability that a given district receives an
engineer Collector, while being unrelated to that district's programme needs.

A potential concern is that the electoral-turnover control---included to account for
politically motivated posting decisions---might itself absorb variation in queue
composition, since governments approaching elections may have stronger incentives to
deploy engineer Collectors strategically. Dropping this control affects neither the
structural coefficient nor the first stage in any material way: the sign-contrast
coefficient changes by less than $0.001$ and the first-stage $F$-statistic moves
from $34.5$ to $29.4$, a reduction that does not alter the inference. The IV estimate
is not an artefact of the election control.

Two features of the result warrant emphasis rather than concealment. First, the
estimate attains significance at the ten percent level, not the five percent level;
the strength of the first stage lends confidence to the point estimate but does not
narrow the confidence interval, which includes zero at conventional levels. The
appropriate reading is that instrumenting recovers an effect of the predicted sign and
plausible magnitude on a well-identified first stage, not that comparative advantage
is established beyond doubt. Second, the single-outcome first stage is weak in the
roads-only subsample ($F=3.2$, $N=453$): with fewer than five hundred observations
the within-district-year variation is too thin to identify the effect from queue
composition alone. This is not a defect of the instrument but a demonstration of why
the pooled design is necessary---stacking the two programme margins is what renders
identification feasible at this sample size. In the full officer-district-year panel the
unstacked instrument is powerful ($F=170$ over all available years; $F=44.7$ restricted
to the post-2001 roads sample period), confirming that the weakness is a consequence of
the roads subsample's limited size rather than a problem with the instrument itself.

%% ------------------------------------------------------------------
\subsection{The Falsification Outcome}
\label{sec:results-falsification}

Proposition~\ref{prop:mix} predicts that education shifts the \emph{composition} of
output while leaving its aggregate approximately unchanged: an engineer Collector
delivers more infrastructure and less welfare, but the total value of public service
output is roughly constant. A measure that aggregates over both programme classes
should therefore display no systematic association with officer type; the two effects
cancel in the aggregate.

Night-time luminosity, which reflects economic activity produced by both infrastructure
(roads facilitating commerce) and welfare (employment income raising household
consumption), is well suited to this test. The engineer coefficient on rural nightlights
is $-0.0007$ ($\text{SE}=0.0007$, $p=0.30$), indistinguishable from zero and stable
across all four fixed-effect specifications examined in Section~\ref{sec:results-ols}.
The infrastructure gain and welfare loss offset in aggregate light output, consistent
with the comparative-advantage mechanism.

This null is consistent with the theory but does not by itself establish it, since an
absence of any effect on either outcome would produce the same result. Its role is
to rule out a specific alternative: the generic-ability account, under which engineers
are more capable administrators in all dimensions, would predict a \emph{positive}
aggregate effect. The near-zero nightlights coefficient rules out that account. The
reason we find a null is not that engineers do nothing; it is, according to the
mechanism, that what they do more of on one margin is offset by what they do less of
on the other.

%% ------------------------------------------------------------------
\subsection{The Tenure Prediction Is Not Supported}
\label{sec:results-tenure}

Proposition~\ref{prop:tenure} predicts that the education effect is strongest early
in the career, when accumulated experience has not yet diluted the educational
component of capacity, and fades as years of service accumulate. The data do not
support this.

The split-sample test (Table~\ref{tab:ideas_summary}, Panel~A, rows 2a--2b) divides
the sample at the median years of service, nine years. Among junior Collectors (at or
below the median, $N=644$) the sign-contrast coefficient is $-0.070$ ($p=0.66$);
among senior Collectors (above the median, $N=612$) it is $+0.459$ ($p=0.12$). The
senior estimate is of the predicted sign and is marginally significant at the fifteen
percent level; the junior estimate is of the wrong sign and is entirely uninformative.
The pattern---larger coefficient for senior DCs, smaller or negative for junior---is
the opposite of Proposition~\ref{prop:tenure}'s prediction, though neither subsample
estimate is significant and both rest on limited samples that are sensitive to
distributional assumptions.

The continuous-interaction specification provides additional resolution without
imposing a binary cut. Interacting the engineer indicator with years of service in
the infrastructure regression and including the same district-plus-state-plus-year
fixed effects, the interaction term is $-0.004$ ($\text{SE}=0.021$, $p=0.85$). The
negative sign is directionally consistent with a fading effect---the engineer
coefficient diminishes by an estimated $0.004$ standard deviations per additional
year of service---but is far from statistically distinguishable from zero, providing
no informative evidence for or against the attenuation story. Partitioning years of
service into terciles confirms the absence of a monotone gradient: the engineer
coefficient is $+0.030$ in the first tercile (least experienced), $+0.004$ in the
second, and $+0.019$ in the third. This is neither the monotone decline predicted by
Proposition~\ref{prop:tenure} nor a pattern that suggests any alternative simple
structural interpretation. The middle tercile is the weakest, which is inconsistent
with both a fading-experience story and a steepening-experience story.

We report this failed prediction rather than suppress it. One interpretation is
admitted by the framework itself: if experience is task-specific rather than
task-general (Appendix~\ref{app:dynamic}), officers accumulate human capital in the
direction their education already favours, so the education-aligned performance gap
widens with tenure rather than narrowing. This would generate a positive interaction
of engineer with years of service, reversing the sign of the prediction. We do not
advance this as an explanation: the tercile pattern is too irregular to be consistent
with monotone task-specific learning, the subsamples are too small to distinguish any
structural story from noise, and the task-general assumption is maintained in the
main theoretical analysis. The honest conclusion is that the tenure channel is not
identified at this sample size, and that the point estimates---to the extent they say
anything---are inconsistent with the baseline prediction rather than supportive of it.

%% ------------------------------------------------------------------
\subsection{Summary}
\label{sec:results-summary}

The pattern across specifications is coherent. The average contrast between engineer
and social-science Collectors is attenuated to insignificance in every OLS
formulation---individual outcome regressions, the pooled signed-outcome regression
($b=+0.037$, $p=0.26$), the explicit sign-contrast interaction ($b=+0.14$, $p=0.21$),
and the log-ratio outcome ($b=+0.06$ to $+0.09$, $p>0.29$)---exactly as the
measurement-error structure of the model predicts. Instrumenting to correct that
attenuation recovers an effect of the predicted sign, roughly nine times larger,
significant at the ten percent level on a first stage well above conventional strength
(KP-$F=34.5$, robust across all control and fixed-effect variations). The aggregate
falsification outcome shows no effect, ruling out a generic-ability account that would
predict a positive aggregate engineer effect. Of the framework's four structural
predictions---attenuation of the OLS contrast, a positive and larger IV estimate,
aggregate output offset, and a fading tenure gradient---three are borne out and one
is not. Section~\ref{sec:counterfactual} takes the instrumented estimate as the basis
for quantifying the output at stake under the seniority assignment the service currently
employs.
